\section{Introduction}
\label{sec:intro}

A typical robot speech pipeline reads one channel from a USB microphone
array and denoises it with a neural model. It passes the result to a
recognizer. Ours reads the processed channel of a ReSpeaker array,
DeepFilterNet3~\cite{schroeter2023deepfilternet3} cleans it with an
attenuation limit of 12\,dB, and Whisper~\cite{radford2023whisper}
transcribes it. The other four channels, the raw microphones, are thrown
away. Two questions follow:
\begin{enumerate}
\item What do those microphones add when the noise comes from a direction?
\item How should the denoiser's only knob, \texttt{atten\_lim\_db}, be set
  when the listener is a recognizer rather than a person?
\end{enumerate}

A first intuition is to switch off the microphone that faces the noise.
It is worth testing because it is how one would reason about directional
microphones. The array's microphones are omnidirectional and 4.5\,cm
apart, so the expectation is that it cannot work. The tool that should
work is a beamformer that combines all four with phase. Whether it works
on this array, in a real room, with noise statistics a robot can actually
obtain, is an empirical question. Project~15, our earlier work on this
array, answered it only in simulation.

\paragraph{Contributions.}
\begin{enumerate}
\item \textbf{A recording protocol with exact scoring.} Chirp-aligned
  speech-only and noise-only takes on a fixed array give exact speech and
  noise components for every linear front end, and real recordings
  underneath. Live takes check the protocol.
\item \textbf{Hardware numbers for nine front ends} on \NScenes{} scenes,
  from the intuitive mute to an RTF-steered MVDR and an adaptive GSC. Each
  uses only the 1.5\,s before the talker speaks.
\item \textbf{A recognition sweep} of DeepFilterNet3's attenuation limit,
  continuous and in the robot's chunked mode, with \NTranscripts{} Whisper
  transcriptions.
\item \textbf{A live comparison with the array's own processed channel.}
\end{enumerate}
